\section*{РЕФЕРАТ}
\addcontentsline{toc}{section}{РЕФЕРАТ}

Расчетно--пояснительная записка 86 с., \totalfigures\ рис., \totaltables\ табл., 27 ист., 3 прил.

Работа посвящена разработке метода оценки кредитоспособности физических лиц с применением алгоритма градиентного бустинга деревьев решений.

\textbf{Ключевые слова:} кредитный скоринг, кредитный риск, машинное обучение, градиентный бустинг, GBDT, классификация, банковские системы.

Рассмотрена задача оценки кредитоспособности заемщиков в банковских системах. Проведен анализ предметной области кредитного скоринга, а также рассмотрены существующие методы, включая традиционные статистические подходы и современные алгоритмы машинного обучения. Выполнен обзор существующих скоринговых систем и проведен сравнительный анализ методов.

Сформулирована цель работы и выполнена формализация постановки задачи с использованием IDEF0-диаграммы.

Метод оценки кредитоспособности разработан на основе алгоритма градиентного бустинга деревьев решений. Описаны его основные этапы: предобработка данных, обучение модели и интерпретация результатов. Определены структуры данных и схема взаимодействия компонентов системы.

Обоснован выбор инструментов программной реализации. Разработано программное обеспечение, реализующее предложенный метод, описаны входные и выходные данные, а также интерфейсы взаимодействия системы.

Проведено тестирование разработанного программного обеспечения и исследование его эффективности. Выполнено сравнение результатов с традиционными методами кредитного скоринга и проанализировано влияние различных факторов на точность предсказаний.